% --------------------------------------------------------------------------SEITENLAYOUT--------------------------------------------------------------------------------------------

% Informationen Über das Dokument, wie z.B. Titel, Autor, Matrikelnr. etc werden in der Datei _Meta.tex definiert und
% können danach global verwendet werden.

\thispagestyle{plain}								%setze Seitenlayout
% @ param \thispagestyle: plain, empty, headings, myheadings.
%-----------------------------------------------------------------------------------------------------------------------------------------------------------------------------------------
%INFO:
%Legt die Art des Seitenformats für eine einzelne Seite fest. Die Voreinstellung plain steht für eine zentrierte Seitennummer am Seitenfuss.
%Durch empty erreicht man Seiten, die keinerlei Kopf oder Fuß enthalten. headings erzeugt eine Kopfzeile aus der Seitennummer und der Überschrift des
%laufenden Abschnitts. myheadings schließlich erlaubt es, den Seitenkopf selbst mit den \markright- und \markboth-Befehlen zu definieren.

%------------------------------------------------------------------------------------------------------------------------------------------------------------------------------------------

% ---------------------------------------------------------------------------SETZE TITELBLATT----------------------------------------------------------------------------------------

\begin{titlepage}									%

\begin{center} 									% setze die Formatierung auf Center
\includegraphics[scale=0.2]{pic/zhaw_soe_d1.png}\\[10ex]		% binde die Grafik ein
\huge{\textbf{\textsc{\titel}}}\\[5ex]						% setze eine Huge-Schrift ein (hier der Titel \titel)
\LARGE{\textbf{\untertitel}}\\[4ex]						% setze Large-Schrift (mit Untertitel)
%\LARGE{\textbf{\art}}\\[8ex]							% setze Large Schrift mit der Art der Arbeit
%\Large{im Fachgebiet \fachgebiet}\\[6ex]					% setze Large Schrift mit dem Fachgebiet
\normalsize										% setze zurück auf Normalschrift	

\begin{tabular}{w{3cm}p{6cm}}\\						% mach eine Tabelle mit den Studenten....
 Programmer:	 & \quad \autorFirst\\[1.2ex]					% ...
 %	 & \quad \autorSecond\\[6ex]						% ...
 %Studienbereich: & \quad \studienbereich\\[1.2ex]				% ...
 %Matrikelnummer: & \quad \matrikelnr\\[1.2ex]				% ...
 %Dozent:         & \quad \erstgutachter\\[1.2ex]					% ...
 %Zweitgutachter:         & \quad \zweitgutachter\\[3ex]			% ...
\end{tabular}									% ...
\\[4ex]										% dies ist der Abstand....
\copyright\ \jahr\\[4ex]								% ...
\end{center}										% Ende der Tabelle

\singlespacing									% einfacher Zeilenabstand
\small											% setze kleine Schrift
\noindent										% Sorgt am Anfang eines Absatzes dafür, daß die erste Zeile nicht eingerückt wird.
\\
\\
\\
\\
%Dieses Werk einschließlich seiner Teile ist \textbf{urheberrechtlich geschützt}. Jede Verwertung ausserhalb der engen Grenzen des Urheberrechtgesetzes ist ohne Zustimmung des Autors unzulässig und strafbar. Das gilt insbesondere für Vervielfältigungen, Übersetzungen, Mikroverfilmungen sowie die Einspeicherung und Verarbeitung in elektronischen Systemen.

\end{titlepage}									%

%---------------------------------------------------------------------------------------------------------------------------------------------------------------------------------------------